\documentclass[12pt]{article}
\usepackage{latexblog}

% ============================================================
% Blog Metadata
% ============================================================
\blogtitle{Burrows-Wheeler变换(Burrows–Wheeler Transform)}
\blogdate{2020-03-04}
\blogtags{刨根问底, 算法}
\bloglang{zh}
\blogsummary{}

\begin{document}
\makeblogtitle

最近听一个医学专业的同学提到了在进行基因分析中用到BWT算法，觉得挺有意思的，正巧赶上这次疫情在家，于是想研究一下这个算法。这个算法的核心思想在于，调整原来的字符串中字符的顺序（而不改变其长度及内容）从而更多的将重复的字符排列到一起，这样有助于其他的压缩算法获得更高的压缩比。这个算法在基因分析中大有用处也就顺理成章了，想想DNA的双链表示大概都是G-T-A-C会有很多这样的字符，那么运用BWT应该可以有比较好的效果。

\section{算法实现}\label{ux7b97ux6cd5ux5b9eux73b0}

考虑一个字符串，要想将相同的字符排列到一起，那么最简单的办法就是，将字符串中的字符进行排序。可是单纯的排序之后，虽然还是那么多字符，但是丢失了一个重要的信息就是字符原来的顺序，而BWT的核心思想就是在于排序并想办法保存字符的顺序信息。

\subsection{编码}\label{ux7f16ux7801}

编码方式如下：

\begin{itemize}
\tightlist
\item
  将 \({\displaystyle \$}\) 作为字符串结尾标记加入到原字符串（记为 \({\displaystyle S}\) )末尾
\item
  将字符串从左到右进行轮换，对于一个长度为 \({\displaystyle N}\) 的字符串，产生 \({\displaystyle N}\) 个新的字符串，记为 \({\displaystyle S_{n}}\)
\item
  将 \({\displaystyle S, S_{1}...S_{n}}\) 进行字典序排序， \({\displaystyle \$}\) 权值最小放在最前面，也即 \({\displaystyle S_{n}}\) 在第一个
\item
  取排序后的所有字符串的最后一个字符，生成一个新的字符串（记为\({\displaystyle S^{"}}\) )，即编码完成
\end{itemize}

以字符串\texttt{banana}举例来说：

\[
S = banana, N = 6\\
\begin{cases}
S_{0} = banana\color{blue}{$}\\
S_{1} = anana\color{blue}{$}b\\
S_{2} = nana\color{blue}{$}ba\\
S_{3} = ana\color{blue}{$}ban\\
S_{4} = na\color{blue}{$}bana\\
S_{5} = a\color{blue}{$}banan\\
S_{6} = \color{blue}{$}banana
\end{cases}
\xrightarrow[\text{字典序排序}]{}
\begin{cases}
S_{6} = \color{blue}{$}banana\\
S_{5} = a\color{blue}{$}banan\\
S_{3} = ana\color{blue}{$}ban\\
S_{1} = anana\color{blue}{$}b\\
S_{0} = banana\color{blue}{$}\\
S_{4} = na\color{blue}{$}bana\\
S_{2} = nana\color{blue}{$}ba
\end{cases}
\xrightarrow[\text{获取最后一个字符}]{}
\begin{cases}
S_{6} = $banan\color{red}{a}\\
S_{5} = a$bana\color{red}{n}\\
S_{3} = ana$ba\color{red}{n}\\
S_{1} = anana$\color{red}{b}\\
S_{0} = banana\color{red}{$}\\
S_{4} = na$ban\color{red}{a}\\
S_{2} = nana$b\color{red}{a}
\end{cases}\\
S^{"} = BWT(banana) = annb$aa
\]

可以看出，转换后的字符串\texttt{annb\$aa}比原来的字符串重复相连的字符的确更多了。实际上\href{http://www.bzip.org/}{bzip}就是应用了BWT结合进行压缩的：

\begin{quote}
bzip2 compression program is based on Burrows--Wheeler algorithm.
\end{quote}

BWT转换后的重复相连字符更多并不绝对，有时候可能转换后的情况反而更糟，比如这个例子：

\[
BWT(appellee) = e$elplepa
\]

反而不如原始字符串了。

\subsection{解码}\label{ux89e3ux7801}

\subsubsection{利用还原矩阵法}\label{ux5229ux7528ux8fd8ux539fux77e9ux9635ux6cd5}

解码的过程分为以下几步：

\begin{itemize}
\tightlist
\item
  根据编码后的字符串 \({\displaystyle S^{"}}\) ，得到还原矩阵
\item
  根据还原矩阵，逐个还原出原来的顺序
\end{itemize}

根据编码的过程我们知道，实际上是这样的对应： \[
\begin{cases}
S_{6} = \color{green}{$}banan\color{red}{a}\\
S_{5} = \color{green}{a}$bana\color{red}{n}\\
S_{3} = \color{green}{a}na$ba\color{red}{n}\\
S_{1} = \color{green}{a}nana$\color{red}{b}\\
S_{0} = \color{green}{b}anana\color{red}{$}\\
S_{4} = \color{green}{n}a$ban\color{red}{a}\\
S_{2} = \color{green}{n}ana$b\color{red}{a}
\end{cases}
\xrightarrow[\text{还原矩阵}]{}
\begin{pmatrix}
$ & a\\
a & n\\
a & n\\
a & b\\
b & $\\
n & a\\
n & a
\end{pmatrix}
\]

得到这个矩阵非常简单，直接将字符串 \({\displaystyle S^{"}}\) 排个序就可以得到：

\[
\begin{cases}
------a\\
------n\\
------n\\
------b\\
------$\\
------a\\
------a
\end{cases}
\xrightarrow[\text{排序}]{}
\begin{cases}
------$\\
------a\\
------a\\
------a\\
------b\\
------n\\
------n
\end{cases}
\xrightarrow[\text{还原矩阵}]{}
\begin{pmatrix}
$ & a\\
a & n\\
a & n\\
a & b\\
b & $\\
n & a\\
n & a
\end{pmatrix}
\]

在这样的一个还原矩阵中，每一个字符对应的就是它最末尾的字符。解码的过程如下：

\begin{itemize}
\tightlist
\item
  从左边列的 \({\displaystyle S}\) 开始，找到对应的字符作为下一个字符 \({\displaystyle C_{n}}\)
\item
  根据 \({\displaystyle C_{n}}\) 这个字符，在左边列找到对应的字符，其对应的字符即 \({\displaystyle C_{n-1}}\)
\item
  以此类推，直到结尾
\item
  如果出现了多个相同的字符，那么就从上到下按顺序找就可以了
\end{itemize}

\[
\begin{pmatrix}
\color{red}{$} & \color{red}{a}\\
a & n\\
a & n\\
a & b\\
b & $\\
n & a\\
n & a
\end{pmatrix}
\xrightarrow[\text{\$a}]{}
\begin{pmatrix}
\color{gray}{$} & \color{gray}{a}\\
\color{red}{a} & \color{red}{n}\\
a & n\\
a & b\\
b & $\\
n & a\\
n & a
\end{pmatrix}
\xrightarrow[\text{\$an}]{}
\begin{pmatrix}
\color{gray}{$} & \color{gray}{a}\\
\color{gray}{a} & \color{gray}{n}\\
a & n\\
a & b\\
b & $\\
\color{red}{n} & \color{red}{a}\\
n & a
\end{pmatrix}
\xrightarrow[\text{\$ana}]{}
\begin{pmatrix}
\color{gray}{$} & \color{gray}{a}\\
\color{gray}{a} & \color{gray}{n}\\
\color{red}{a} & \color{red}{n}\\
a & b\\
b & $\\
\color{gray}{n} & \color{gray}{a}\\
n & a
\end{pmatrix}
\xrightarrow[\text{\$anan}]{}
\begin{pmatrix}
\color{gray}{$} & \color{gray}{a}\\
\color{gray}{a} & \color{gray}{n}\\
\color{gray}{a} & \color{gray}{n}\\
a & b\\
b & $\\
\color{gray}{n} & \color{gray}{a}\\
\color{red}{n} & \color{red}{a}\\
\end{pmatrix}
\xrightarrow[\text{\$anana}]{}
\begin{pmatrix}
\color{gray}{$} & \color{gray}{a}\\
\color{gray}{a} & \color{gray}{n}\\
\color{gray}{a} & \color{gray}{n}\\
\color{red}{a} & \color{red}{b}\\
b & $\\
\color{gray}{n} & \color{gray}{a}\\
\color{gray}{n} & \color{gray}{a}\\
\end{pmatrix}
\xrightarrow[\text{\$ananab}]{}
\begin{pmatrix}
\color{gray}{$} & \color{gray}{a}\\
\color{gray}{a} & \color{gray}{n}\\
\color{gray}{a} & \color{gray}{n}\\
\color{gray}{a} & \color{gray}{b}\\
\color{red}{b} & \color{red}{$}\\
\color{gray}{n} & \color{gray}{a}\\
\color{gray}{n} & \color{gray}{a}\\
\end{pmatrix}\\
S^{'} = \$ananab\\
S = reverse(S^{'}) \\= banana\$
\]

\subsubsection{变种}\label{ux53d8ux79cd}

另一种方式可能更清晰，但实质上是一回事，只是做法看着不一样。在上述构建还原矩阵的过程中，我们实际已知的是最后一列的数据，那么，如果我们想办法把其他的列都构建出来，就可以得到原来的字符串了。

\[
\begin{cases}
------a\\
------n\\
------n\\
------b\\
------$\\
------a\\
------a
\end{cases}
\xrightarrow[\text{想办法变成这样}]{}
\begin{cases}
S_{6} = \color{green}{$}banan\color{red}{a}\\
S_{5} = \color{green}{a}$bana\color{red}{n}\\
S_{3} = \color{green}{a}na$ba\color{red}{n}\\
S_{1} = \color{green}{a}nana$\color{red}{b}\\
S_{0} = \color{green}{b}anana\color{red}{$}\\
S_{4} = \color{green}{n}a$ban\color{red}{a}\\
S_{2} = \color{green}{n}ana$b\color{red}{a}
\end{cases}
\xrightarrow[\text{然后拿到S6就可以了}]{}
S_{6} = \color{green}{$}banan\color{red}{a}
\]

过程是这样的： \[
\begin{cases}
------a\\
------n\\
------n\\
------b\\
------$\\
------a\\
------a
\end{cases}
\xrightarrow[\text{将最后一列排序，作为第一列}]{}
\begin{cases}
$-----a\\
a-----n\\
a-----n\\
a-----b\\
b-----$\\
n-----a\\
n-----a
\end{cases}
\]

得到这个之后，从又到左得到\texttt{a\$,na,na,ba,\$b,\ an,\ an}，再将其排序作为第二列, 以此类推：

\[
\begin{cases}
------a\\
------n\\
------n\\
------b\\
------$\\
------a\\
------a
\end{cases}
\rightarrow
\begin{cases}
$-----a\\
a-----n\\
a-----n\\
a-----b\\
b-----$\\
n-----a\\
n-----a
\end{cases}
\rightarrow
\begin{cases}
a$\\
na\\
na\\
ba\\
$b\\
an\\
an
\end{cases}
\xrightarrow[\text{排序}]{}
\begin{cases}
$\color{green}{b}\\
a\color{green}{$}\\
a\color{green}{n}\\
a\color{green}{n}\\
n\color{green}{a}\\
n\color{green}{a}\\
b\color{green}{a}
\end{cases}
\xrightarrow[\text{得到第三列}]{}
\begin{cases}
$b----a\\
a$----n\\
an----n\\
an----b\\
ba----$\\
na----a\\
na----a
\end{cases}\\
\rightarrow
\begin{cases}
a$b\\
na$\\
nan\\
ban\\
$ba\\
ana\\
ana
\end{cases}
\rightarrow
\begin{cases}
$b\color{green}{a}\\
a$\color{green}{b}\\
an\color{green}{a}\\
an\color{green}{a}\\
ba\color{green}{n}\\
na\color{green}{$}\\
na\color{green}{n}
\end{cases}
\xrightarrow[\text{得到第三列}]{}
\begin{cases}
$ba---a\\
a$b---n\\
ana---n\\
ana---b\\
ban---$\\
na$---a\\
nan---a
\end{cases}
\xrightarrow[\text{如此反复，最终得到全部列}]{}
\]

\section{代码实现}\label{ux4ee3ux7801ux5b9eux73b0}

看似复杂的操作，没想到用python可以写的如此简单，不过不见得一看就懂\ldots{}

\begin{Shaded}
\begin{Highlighting}[]
\NormalTok{EOL }\OperatorTok{=} \StringTok{\textquotesingle{}$\textquotesingle{}}

\KeywordTok{def}\NormalTok{ encode(source):}
\NormalTok{    source }\OperatorTok{=}\NormalTok{ source }\OperatorTok{+}\NormalTok{ EOL}
\NormalTok{    table }\OperatorTok{=}\NormalTok{ [source[i:] }\OperatorTok{+}\NormalTok{ source[:i] }\ControlFlowTok{for}\NormalTok{ i }\KeywordTok{in} \BuiltInTok{range}\NormalTok{(}\BuiltInTok{len}\NormalTok{(source))]}
\NormalTok{    table.sort()}

    \ControlFlowTok{return} \StringTok{\textquotesingle{}\textquotesingle{}}\NormalTok{.join([row[}\OperatorTok{{-}}\DecValTok{1}\NormalTok{] }\ControlFlowTok{for}\NormalTok{ row }\KeywordTok{in}\NormalTok{ table])}

\KeywordTok{def}\NormalTok{ decode(encoded):}
\NormalTok{    length }\OperatorTok{=} \BuiltInTok{len}\NormalTok{(encoded)}
\NormalTok{    table }\OperatorTok{=}\NormalTok{ [}\StringTok{\textquotesingle{}\textquotesingle{}}\NormalTok{] }\OperatorTok{*}\NormalTok{ length}

    \ControlFlowTok{for}\NormalTok{ i }\KeywordTok{in} \BuiltInTok{range}\NormalTok{(length):}
\NormalTok{        table }\OperatorTok{=} \BuiltInTok{sorted}\NormalTok{([encoded[m] }\OperatorTok{+}\NormalTok{ table[m] }\ControlFlowTok{for}\NormalTok{ m }\KeywordTok{in} \BuiltInTok{range}\NormalTok{(length)])}
        \BuiltInTok{print}\NormalTok{(table)}
\NormalTok{    s }\OperatorTok{=}\NormalTok{ [row }\ControlFlowTok{for}\NormalTok{ row }\KeywordTok{in}\NormalTok{ table }\ControlFlowTok{if}\NormalTok{ row.endswith(EOL)][}\DecValTok{0}\NormalTok{]}
    \ControlFlowTok{return}\NormalTok{ s.rstrip(EOL)}
\end{Highlighting}
\end{Shaded}

解码的这个循环不大好理解，打出来一看就懂了：

\begin{verbatim}
['$', 'a', 'a', 'a', 'b', 'n', 'n']
['$b', 'a$', 'an', 'an', 'ba', 'na', 'na']
['$ba', 'a$b', 'ana', 'ana', 'ban', 'na$', 'nan']
['$ban', 'a$ba', 'ana$', 'anan', 'bana', 'na$b', 'nana']
['$bana', 'a$ban', 'ana$b', 'anana', 'banan', 'na$ba', 'nana$']
['$banan', 'a$bana', 'ana$ba', 'anana$', 'banana', 'na$ban', 'nana$b']
['$banana', 'a$banan', 'ana$ban', 'anana$b', 'banana$', 'na$bana', 'nana$ba']
\end{verbatim}

\begin{itemize}
\tightlist
\item
  \href{https://zh.wikipedia.org/wiki/Burrows-Wheeler\%E5\%8F\%98\%E6\%8D\%A2}{维基百科-Burrows-Wheeler变换}
\item
  \href{https://www.cs.cmu.edu/~ckingsf/bioinfo-lectures/bwt.pdf}{BWT}
\end{itemize}


\end{document}
